\documentclass[11pt]{article}
\usepackage{fontspec}
\usepackage[french]{babel}
\usepackage{amsmath}
\usepackage{amssymb}
\usepackage{amsthm}
\usepackage{mathrsfs}
\usepackage{tikz-cd}
\usepackage{enumitem}

\newcommand{\printedpage}[1]{\marginpar{\footnotesize p.~#1}}
\newcommand{\msp}[1]{\marginpar{\footnotesize [#1]}}
\newcommand{\cercle}[1]{\tikz[baseline=(c.base)]\node[draw,circle,inner sep=1pt](c){#1};}

\begin{document}

\printedpage{59}
% head: A. Grothendieck : Dérivateurs, Ch. XVII, Édition des Universités de Montpellier II et Paris VII
% foot: Version 9 février 2019

Or la relation en question $h_0\tilde{i} = \tilde{h}_0$ équivaut à l'ensemble des deux relations
\[
\begin{aligned}
(h_0\tilde{i})r &= h_0(\tilde{i}r)\\
&= h_0i_0\\
&\overset{?}{=} \alpha_0\\
&= \tilde{h}_0r,
\end{aligned}
\]
équivalent à $h_0i_0 = \alpha_0$,
\[
\begin{aligned}
(h_0\tilde{i})s &= h_0(\tilde{i}s)\\
&= h_0\lambda_B\\
&\overset{?}{=} \lambda_Xh'_{I'}\\
&= \tilde{h}_0s,
\end{aligned}
\]
équivalent à $h_0\lambda_B = \lambda_Xh'_{I'}$, c'est OK.

Ceci posé, dans le carré (6.3.14), les hypothèses $i \in \Phi(I)$, $p \in \Psi\langle I\rangle$ impliquent $\tilde{i}_0 \in \Phi$, $p_o \in \Psi$. Donc la flèche cherchée $h_0$ existe, q.e.d. % [sic] p_o -> p_0, index printed as letter o, elsewhere index 0 (h_0, i_0)

Cela achève la démonstration du corollaire 6.3, et par là aussi, celle de la dernière partie, d), de 6.1. On est heureux\,!

\msp{page 117}

\textbf{Proposition 6.4.}\footnote[34]{Sous réserve d'existence des $\varinjlim$ et des $\varprojlim$ pertinentes dans $\mathcal{M}$. On peut supposer pour simplifier $\mathcal{M}$ stable par $\varinjlim$ et $\varprojlim$ .} % mark printed as superscript (34)
\textit{Soit $\mathcal{M}$ une catégorie, et $\Phi \longleftrightarrow \Psi$ un couple Q-clos en Q-dualité ($\Phi = \Psi^*$, $\Psi = \Phi_*$). Soit $u : I \longrightarrow I'$ une flèche dans} Cat\textit{, d'où une suite de foncteurs adjoints}
\[
\begin{tikzcd}[column sep=large]
\mathcal{M}(I) \arrow[r, shift left=3, "u_!"] \arrow[r, shift right=3, "u_*"'] & \mathcal{M}(I') \arrow[l, "u^*"']
\end{tikzcd}\ .
\]
\textit{On définit $\Phi(I)$, $\Psi\langle I\rangle$, $\Phi(I')$, $\Psi\langle I'\rangle$ comme dans} 6.2.5. \textit{On a alors :}

\cercle{a}
\[
\left\{
\begin{aligned}
u_!(\Phi(I)) &\ \subset\ \Phi(I')\ ,\\
u^*(\Psi\langle I'\rangle) &\ \subset\ \Psi\langle I\rangle\ .
\end{aligned}
\right.
\tag*{(6.4.1)}
\]

\cercle{b} \textit{On a l'équivalence}
\[
u^*(\Phi(I')) \subset \Phi(I) \qquad\Longleftrightarrow\qquad u_*(\Psi\langle I\rangle) \subset \Psi\langle I'\rangle\ .
\tag*{(6.4.2)}
\]

\cercle{c} \textit{Les relations équivalentes} (6.4.2) \textit{ont lieu dans chacun des trois cas suivants :}
\begin{enumerate}[label={}, leftmargin=2.5em]
\item 1°) \textit{$u$ est coétale.}
\item 2°) \textit{$I$ est discrète.}
\item 3°) \textit{Pour tout $i' \in \operatorname{Ob} I'$, les composantes connexes de $I/i'$ admettent un objet final.}
\end{enumerate}

\newpage
\printedpage{60}
% head: A. Grothendieck : Dérivateurs, Ch. XVII, Édition des Universités de Montpellier II et Paris VII
% foot: Version 9 février 2019

\textbf{Corollaire 6.5.} \textit{Soit $(W,C,F)$ un triple de Quillen clos spécial dans $\mathcal{M}$, et définissons, pour toute $I$ dans} Cat\textit{, $C(I)$, $TC(I)$ et $F(I)$, $TF(I)$ comme dans le théorème} 6.1. a)\textit{. On a alors}

\msp{page 118}

\textit{les équivalences}
\[
\Bigl(u^*(C(I'))\ \subset\ C(I)\Bigr) \quad\Longleftrightarrow\quad \Bigl(u_*(TF)\ \subset\ TF'\Bigr)\ ,
\tag*{(6.5.1)}
\]
\[
\Bigl(u^*(TC(I'))\ \subset\ TC(I)\Bigr) \quad\Longleftrightarrow\quad \Bigl(u_*(F)\ \subset\ F'\Bigr)\ .
\tag*{(6.5.2)}
\]
\textit{Les conditions équivalentes} (6.5.1) \textit{ainsi que} (6.5.2) \textit{sont satisfaites quand on est dans l'un des cas} 1°) 2°) 3°) \textit{de} 6.4. c).

\textsc{Démonstration} de 6.4. Les assertions \cercle{a} , \cercle{b} sont $\pm$ tautologiques, par adjonction, comparer XIII, prop. 2, et cor., page 46. Pour prouver c), il suffit de traiter le cas 3°), car 1° $\Longrightarrow$ 3°, 2° $\Longrightarrow$ 3°. Mais le calcul de
\[
u_*(X)(i')\ \simeq\ \varprojlim_{i\in I/i'} X(i)
\]
montre que si les composantes connexes $(I/i')_\alpha$ de $I/i'$ ont chacune un objet final $e_\alpha$, alors
\[
u_*(X)(i')\ \simeq\ \prod_\alpha X(e_\alpha)\ .
\]
Or $\Psi$, étant clos à droite, est stable par produits quelconques, donc $X \in \Psi\langle I\rangle$ implique que les $u_*(X)(i')$ sont dans $\Psi$, donc $u_*$ est dans $\Psi\langle I'\rangle$, q.e.d.

\end{document}
